\pagenumbering{arabic}
\section{INTRODUCCIÓN}

La digitalización del mercado de boletería ha facilitado la compra y la transferencia de entradas, pero
también ha 
incrementado los riesgos asociados a la reventa fraudulenta: duplicidad de códigos de acceso, venta
concurrente del 
mismo activo a varios compradores y opacidad del historial de propiedad. Estos vectores afectan la confianza
del 
comprador, la liquidación efectiva entre los actores del mercado y la trazabilidad operacional del
organizador.

Este documento presenta el diseño, la implementación y la validación de \textbf{SecureTicket}, una solución
B2B de 
reventa de entradas basada en contratos inteligentes Soroban sobre la red Stellar. La propuesta materializa
cada 
boleto como un \textit{token} no fungible (NFT) con versionado, aplica un patrón \textit{burn/remint} en cada
reventa 
para preservar la unicidad y ejecuta la liquidación de pagos y comisiones de forma atómica dentro de la misma
invocación del contrato.

El documento se organiza en nueve secciones. La sección~2 presenta la descripción general del proyecto,
incluyendo 
problema, propuesta de solución, objetivos y entregables. La sección~3 expone el contexto técnico y el
análisis de 
soluciones comparables. La sección~4 detalla los requerimientos, restricciones y la especificación funcional
del sistema. 
La sección~5 desarrolla la arquitectura y el diseño detallado de los contratos y los nodos de la solución. La
sección~6 
describe el proceso de implementación, despliegue e integración continua. La sección~7 consolida los
resultados 
experimentales sobre la red de prueba (\textit{testnet}). La sección~8 analiza el impacto del trabajo y
presenta las 
conclusiones y el trabajo futuro. La sección~9 lista las referencias bibliográficas. Finalmente, los anexos
referencian 
el Plan de Gestión del Proyecto, la Especificación de Requisitos de Software y el documento de Diseño y
Propuesta del 
Trabajo de Grado, entregados como artefactos independientes.
